% Slides for the mesh-study meeting. Same pictures and tables as mesh_study.tex.
% Build with:  sh DLR/meshStudy/report/build.sh   (Overleaf: pdfLaTeX)
\documentclass[aspectratio=169,10pt]{beamer}
\usetheme{default}
\usecolortheme{dove}
\setbeamertemplate{navigation symbols}{}
\setbeamertemplate{footline}[frame number]
\setbeamerfont{caption}{size=\scriptsize}
\usepackage[T1]{fontenc}
\usepackage[utf8]{inputenc}
\usepackage{lmodern}
\usepackage{amsmath}
\usepackage{booktabs}
\graphicspath{{fig/}}

% \pic{width}{file}: the picture, or a placeholder box if it has not been generated yet
\newcommand{\pic}[2]{\IfFileExists{fig/#2}{\includegraphics[width=#1]{#2}}{%
  \fbox{\parbox[c][2.2cm][c]{\dimexpr#1-2\fboxsep-2\fboxrule}{\centering\scriptsize
  \texttt{\detokenize{#2}}\\not generated yet}}}}
% \tab{file}: a generated tabular, scaled to the slide width if needed
\newcommand{\tab}[2]{\IfFileExists{#1}{\resizebox{#2}{!}{\input{#1}}}{\textit{#1 missing: run make\_tex.py}}}

\IfFileExists{summary.tex}{% generated by make_tex.py -- do not edit, rerun build.sh instead
\newcommand{\StudyStatus}{Status of the runs when this document was built (02:21 on 17 Nov 2026): coarse: finished; medium: meshing; fine: not started.}
\newcommand{\CellsCoarse}{1\,190\,168}
\newcommand{\CellsMedium}{--}
\newcommand{\CellsFine}{--}
\newcommand{\FindingsQuality}{\begin{itemize}\item coarse mesh (1\,190\,168 cells): fails 1 checkMesh check: max skewness 8.61 on 5 faces (limit 4). The skewed faces are all outside the refinement box; the nearest is 81\,m from a probe.\end{itemize}}
\newcommand{\FindingsConvergence}{\begin{itemize}\item coarse: initial residuals after 1000 iterations between $2.4\times10^{-9}$ ($U_x$) and $3.0\times10^{-6}$ ($p$); continuity error (sum local) $3.1\times10^{-11}$; the probes moved by less than 0.01\,\% over the last 100 iterations.\end{itemize}}
\newcommand{\FindingsIndependence}{\textit{At least two finished runs are needed to compare meshes; not available yet.}}
\newcommand{\SkewSentence}{The coarse mesh fails the skewness check: the maximum skewness is 8.61 on 5 of its 3\,708\,464 faces, against a usual limit of 4.}
\newcommand{\NFinished}{1}
}{%
  \newcommand{\StudyStatus}{summary.tex is missing: run make\_tex.py.}%
  \newcommand{\CellsCoarse}{--}\newcommand{\CellsMedium}{--}\newcommand{\CellsFine}{--}%
  \newcommand{\FindingsQuality}{}\newcommand{\FindingsConvergence}{}\newcommand{\FindingsIndependence}{}%
  \newcommand{\SkewSentence}{}\newcommand{\NFinished}{0}}

\title{Mesh-independence study\\of the DLR Oberpfaffenhofen wind model}
\author{Henry Delacroix}
\institute{Politecnico di Milano, Multidisciplinary Project SWARM-LINK}
\date{\today}

\begin{document}

\begin{frame}
\titlepage
\vspace{-1em}
\begin{center}\scriptsize\textit{\StudyStatus}\end{center}
\end{frame}

% ---------------------------------------------------------------------------------------------
\begin{frame}{Goal and set-up}
\begin{columns}[T]
\begin{column}{0.40\textwidth}\small
\textbf{Goal:} refine the campus mesh to drone scale ($\le1$\,m cells) around a few buildings and show
that the wind no longer depends on the mesh.

\medskip
\begin{itemize}\setlength{\itemsep}{2pt}
  \item Steady RANS $k$--$\varepsilon$, \texttt{simpleFoam} (OpenFOAM v2306)
  \item West wind, 5\,m/s at 10\,m, log-law inflow
  \item 1000 iterations per mesh, no turbine
  \item 3 meshes: \textbf{only} the cell size in the red box changes: 4 / 2 / 1\,m (ratio 2)
  \item Box covers the weather station $\pm50$\,m, for later validation
  \item 19 probes on 4 vertical lines
\end{itemize}
\end{column}
\begin{column}{0.58\textwidth}
\centering\pic{\linewidth}{box_topview.png}
\end{column}
\end{columns}
\end{frame}

% ---------------------------------------------------------------------------------------------
\begin{frame}{The three meshes: horizontal cut at $z=24$\,m (about 7\,m above ground)}
\centering
\begin{minipage}[t]{0.325\textwidth}\centering\pic{\linewidth}{mesh_slice_coarse.png}\\\scriptsize coarse: 4\,m, \CellsCoarse{} cells\end{minipage}\hfill
\begin{minipage}[t]{0.325\textwidth}\centering\pic{\linewidth}{mesh_slice_medium.png}\\\scriptsize medium: 2\,m, \CellsMedium{} cells\end{minipage}\hfill
\begin{minipage}[t]{0.325\textwidth}\centering\pic{\linewidth}{mesh_slice_fine.png}\\\scriptsize fine: 1\,m, \CellsFine{} cells\end{minipage}

\medskip
\scriptsize Each polygon is one cell; grey = buildings; red dashed = refinement box.
\end{frame}

% ---------------------------------------------------------------------------------------------
\begin{frame}{Fine mesh along the vertical (cut at $y=206$\,m, through the station)}
\centering\pic{\textwidth}{mesh_vertical_fine.png}

\medskip
\small 1\,m cells inside the red box (up to $z=60$\,m), then coarser cells (2 and 4\,m) above and around it.
Stars: the station probes, 15--40\,m above ground. Blue dotted: height of the horizontal cut.
\end{frame}

% ---------------------------------------------------------------------------------------------
\begin{frame}{Mesh quality and convergence}
\begin{columns}[T]
\begin{column}{0.49\textwidth}\centering\scriptsize\textbf{Mesh quality} (bold: outside the usual limit)\\[2pt]
\tab{tab_quality.tex}{\linewidth}
\end{column}
\begin{column}{0.49\textwidth}\centering\scriptsize\textbf{Convergence} (last iteration)\\[2pt]
\tab{tab_convergence.tex}{\linewidth}
\end{column}
\end{columns}

\medskip
\centering\pic{0.72\textwidth}{residuals.png}
\end{frame}

% ---------------------------------------------------------------------------------------------
\begin{frame}{Mesh independence: wind speed at the probes}
\centering
\tab{tab_independence.tex}{0.88\textwidth}

\medskip
\begin{minipage}{0.88\textwidth}\scriptsize\FindingsIndependence\end{minipage}
\end{frame}

% ---------------------------------------------------------------------------------------------
\begin{frame}{Differences along the vertical, and the wind on the cut at $z=24$\,m}
\centering
\pic{0.8\textwidth}{mesh_study_profiles.png}

\smallskip
\begin{minipage}[t]{0.28\textwidth}\centering\pic{\linewidth}{wind_slice_coarse.png}\end{minipage}\hfill
\begin{minipage}[t]{0.28\textwidth}\centering\pic{\linewidth}{wind_slice_medium.png}\end{minipage}\hfill
\begin{minipage}[t]{0.28\textwidth}\centering\pic{\linewidth}{wind_slice_fine.png}\end{minipage}
\end{frame}

% ---------------------------------------------------------------------------------------------
\begin{frame}{Caveats and next steps}
\begin{columns}[T]
\begin{column}{0.5\textwidth}\small
\textbf{Caveats}
\begin{itemize}\setlength{\itemsep}{2pt}
  \item \SkewSentence
  \item Station position known to $\pm50$\,m only; mast height unknown
  \item RANS: mean wind only, no gusts
  \item First-order upwind convection (from the tutorial)
  \item GCI is an indicator: snappyHexMesh is not a perfectly systematic refinement
  \item One wind direction, placeholder inflow
\end{itemize}
\end{column}
\begin{column}{0.5\textwidth}\small
\textbf{Next steps}
\begin{enumerate}\setlength{\itemsep}{2pt}
  \item Choose the production mesh from the table
  \item Add the wind turbine (actuator disk)
  \item Compare its wake with published data
  \item Validate with the campus weather station (raw data and mast details from DLR)
  \item Prevailing SW wind, real terrain (DGM1)
\end{enumerate}
\end{column}
\end{columns}
\end{frame}

\end{document}
